\documentclass[11pt]{article}

\usepackage{fullpage}
\usepackage{amsmath}
\usepackage{amssymb}
\usepackage{amsfonts}
\usepackage{amsthm}
\usepackage{mathtools}
\usepackage{enumitem}
\usepackage[colorlinks,linkcolor=magenta,citecolor=blue]{hyperref}

\newcommand{\Rbar}{\bar R}
\newcommand{\Om}{\Omega}
\newcommand{\E}{\mathbb E}
\newcommand{\Var}{\operatorname{Var}}
\newcommand{\tr}{\operatorname{tr}}
\newcommand{\diag}{\operatorname{diag}}
\newcommand{\Unif}{\operatorname{Unif}}
\newcommand{\IG}{\operatorname{IG}}
\newcommand{\Gam}{\operatorname{Gamma}}
\newcommand{\Normal}{\mathcal N}

\newtheorem{lemma}{Lemma}

\title{Better approximation to the susie-love model}
\author{}
\date{}

\begin{document}
\maketitle


\section{Model discussion}
Last time we considered the susie-relax model, where the uncertainty of $R$ only affects the mean, not the variance:
\begin{align}\label{eq:model-susie-relax} 
\begin{cases}
x | R, \beta, \gamma \sim \mathcal{N}(\beta R_{\gamma}, \bar R / N)\\
R \sim \mathcal{IW}(\nu_0 \bar R, \nu_0 + J + 1).
\end{cases}
\end{align}
which leads to the conditional distribution:
\begin{equation}\label{eq:conditional-dist-susie-relax}
p(x_{-j}\mid x_j, \beta, \gamma=1) = 
\Normal\left(
x_j \bar R_{j, -j} ,\,
\left(\frac{1}{N}+ \dfrac{\beta^2}{\nu_0 + 3}\right) \bar S_j
\right),
\end{equation}
where $\bar S_j$ is the Schur complement of $\bar R_{jj}$ in $\bar R$.
Because this depends on $\beta$, this model does not satisfy the ``irrelevance of null SNPs'' property. 
Is there another model that arises from the conditional distribution:
\begin{equation}\label{eq:conditional-dist-Matthew}
p(x_{-j}\mid x_j, \beta, \gamma=j) = 
\Normal\left(
x_j \bar R_{j, -j} ,\,
\left(\frac{1}{N}+ \dfrac{x_j^2}{\nu_0 + 3}\right)S_j
\right)?
\end{equation}

\section{Some new models with better approximation to susie-love}
Recall our original SER model from `susie-love':
\begin{align}\label{eq:likelihood-x-R-mean-var} 
\begin{cases}
x | R, \beta, \gamma \sim \mathcal{N}(\beta R_{\gamma}, R / N)\\
R \sim \mathcal{IW}(\nu_0 \bar R, \nu_0 + J + 1).
\end{cases}
\end{align}
In the following, we show that it satisfies ``irrelevance of the null SNPs'' and, in fact, leads to conditional distribution~\eqref{eq:conditional-dist-Matthew}. The novolty here is what we learned from susie-relax last week: To represent $R$ in terms of blocks.

\subsection{Model A (equivalent to the original SER-love model)} Let's write down $p(x | R, \beta, \gamma = 1)$, other cases of $\gamma$ are similar. Let 
$$d_1 = R_{11} \in \mathbb{R}_{+}, \quad z_1 = \dfrac{R_{-1, 1}}{d_1} \in \mathbb{R}^{J-1},$$
and 
$$S_1 = R_{-1, -1} - d_1 z_1 z_1^\top $$
be the Schur's complement of $R_{11}$ in $R$. Hence, the latent $R$ is represented as
\begin{equation}\label{eq:R-block}
R = \begin{pmatrix}
d_1 & d_1 z_1^\top \\
d_1 z_1 & S_1 + d_1 z_1 z_1^\top
\end{pmatrix}.
\end{equation}
Similarly, denote
$$\bar S_1 = \bar R_{-1, -1} - \bar R_{-1, 1} \bar R_{-1, 1}^\top $$
be the Schur's complement of $\bar R_{11}$ in $\bar R$, and $\mu_1 =  \bar R_{-1, 1}$, we can write
\begin{equation}\label{eq:R-bar-block}
\bar R = \begin{pmatrix}
1 &\mu_1^\top \\
\mu_1 & \bar S_1 + \mu_1 \mu_1^\top
\end{pmatrix}.
\end{equation}
Properties of IW distribution (together with the fact that $\bar R_{11} = 1$)
imply:
\begin{equation}
\begin{cases}
d_1 \sim \IG\left(\dfrac{\nu_0 + 2}{2}, \dfrac{\nu_0}{2}\right) \\
z_1 | S_1 \sim \mathcal{N}\left(\bar R_{-1, 1}, \dfrac{S_1}{\nu_0}\right)\\
S_1 \sim \mathcal{IW}(\nu_0 \bar S_1, \nu_0 + J + 1)
\end{cases}
\end{equation}
The original SER version of Susie-love (with $\gamma = 1$; similar for other $\gamma$) becomes:
\begin{align}\label{eq:SER-model-A} 
\begin{cases}
x | \beta, \gamma = 1, z_1, S_1 \sim \mathcal{N}\left(\beta 
\begin{pmatrix}
d_1  \\
d_1 z_1
\end{pmatrix}, 
\dfrac{1}{N} \begin{pmatrix}
d_1 & d_1 z_1^\top \\
d_1 z_1 & S_1 + d_1 z_1 z_1^\top
\end{pmatrix}\right)\\
d_1 \sim \IG\left(\dfrac{\nu_0 + 2}{2}, \dfrac{\nu_0}{2}\right) \\
z_1 | S_1 \sim \mathcal{N}\left(\mu_1, \dfrac{S_1}{\nu_0}\right)\\
S_1 \sim \mathcal{IW}(\nu_0 \bar S_1, \nu_0 + J + 1)\\
\beta \sim \mathcal{N}(0, \sigma_0^2)
\end{cases}
\end{align}
Write
\(
x=\begin{pmatrix}x_1\\x_{-1}\end{pmatrix}.
\)
From the block normal distribution in~\eqref{eq:SER-model-A}, we have
\[
p(x\mid \gamma=1,\beta,d_1,z_1,S_1)
=
\Normal\left(x_1\mid \beta d_1,\frac{d_1}{N}\right)
\Normal_{J-1}\left(x_{-1}\mid x_1z_1,\frac{S_1}{N}\right).
\]
The key simplification is that, \textbf{after conditioning on $x_1$, the conditional
mean of $x_{-1}$ is $x_1z_1$ and no longer contains $\beta$ or $d_1$} (relevant to our earlier discussion). Now we proceed to show that this satisfies the ``irrelevance of null SNPs'' property:
\paragraph{1. Marginalize out $z_1 | S_1$} leads to:
\[
p(x\mid \gamma=1,\beta,d_1,S_1)
=
\Normal\left(x_1\mid \beta d_1,\frac{d_1}{N}\right)
\Normal_{J-1}\left(
x_{-1}\mid x_1\mu_1,
\left(\frac{1}{N}+\frac{x_1^2}{\nu_0}\right)S_1
\right).
\]
What I found so interesting is that the conditional distribution $p(x_{-1}\mid x_1,S_1,\gamma=1)$ resembles what was suggested by Matthew last time (cf.~\eqref{eq:conditional-dist-Matthew}).

\paragraph{2. Marginalize over $S_1$.}
Thus
\[
p(x\mid \gamma=1,\beta,d_1)
=
\Normal\left(x_1\mid \beta d_1,\frac{d_1}{N}\right)
t_{J-1,\nu_0+3}\left(
x_{-1}\mid x_1\mu_1,\,
\frac{\nu_0}{\nu_0+3}\left(\frac{1}{N}+\frac{x_1^2}{\nu_0}\right)\bar S_1
\right),
\]
where the second argument is the multivariate $t$ scale matrix and is denoted by $T_1(x)$:
\[
T_1(x) \propto |\bar S_1|^{-1/2}
\left(\frac{1}{N}+\frac{x_1^2}{\nu_0}\right)^{-(J-1)/2}
\left[
1+
\frac{
(x_{-1}-x_1\mu_1)^\top(\nu_0\bar S_1)^{-1}
(x_{-1}-x_1\mu_1)
}{\frac{1}{N}+\frac{x_1^2}{\nu_0}}
\right]^{-(\nu_0+J+2)/2}.
\]
Using the block inverse formula (notice the block form of $\bar R$~\eqref{eq:R-bar-block}),
\[
x^\top\bar R^{-1}x
=
x_1^2+
(x_{-1}-x_1\mu_1)^\top\bar S_1^{-1}
(x_{-1}-x_1\mu_1),
\]
so
\[
1+
\frac{
(x_{-1}-x_1\mu_1)^\top(\nu_0\bar S_1)^{-1}
(x_{-1}-x_1\mu_1)
}{\frac{1}{N}+\frac{x_1^2}{\nu_0}}
=
\frac{x^\top\bar R^{-1}x+\nu_0/N}{x_1^2+\nu_0/N}.
\]
Therefore
\begin{equation}\label{eq:T_1}
T_1(x)
\propto
|\bar S_1|^{-1/2}
\left(x_1^2+\frac{\nu_0}{N}\right)^{(\nu_0+3)/2}
\left(x^\top\bar R^{-1}x+\frac{\nu_0}{N}\right)^{-(\nu_0+J+2)/2}.
\end{equation}
When comparing different values of $\gamma$, \textbf{the final factor is common across
SNPs.} Also, since $\bar R_{11}=1$, $|\bar R|=|\bar S_1|$; analogously,
$|\bar S_j|=|\bar R|$ for any SNP $j$ with $\bar R_{jj}=1$.



\paragraph{3. Marginalize over $\beta$ and $d_1$} finally leads to 
\[
\boxed{
p(x\mid \gamma=1)
=
T_1(x)
\int_0^\infty
\Normal\left(x_1\mid 0,\frac{d_1}{N}+d_1^2\sigma_0^2\right)
d\IG\left(d_1;\frac{\nu_0+2}{2},\frac{\nu_0}{2}\right)
}.
\]
This is a one-dimensional integral over $d_1>0$ which does not have closed-form expression. Cancel out the first (determinant term) and the last (quadratic term) in $T_1(x)$ because they are common for all SNPs, we have
\begin{equation}\label{eq:PIP-model-A}
\boxed{
p(\gamma=j\mid x)
\propto
\left(1+\frac{Nx_j^2}{\nu_0}\right)^{(\nu_0+3)/2}
\int_0^\infty
\Normal\left(x_j\mid 0,\frac{d_j}{N}+d_j^2\sigma_0^2\right)\,d \IG(d_j)
}.
\end{equation}
\textbf{Conclusion.} The resulting PIP weight depends on SNP $j$
only through $x_j$.

\subsection{Model B} 
Simplify $d_1 \equiv 1$ to avoid intractable integral (and we know the diagonal of the in-sample LD matrix is 1 anyway), we have the following model (model B):
\begin{align}\label{eq:SER-model-B} 
\begin{cases}
x | \beta, \gamma = 1, z_1, S_1 \sim \mathcal{N}\left(\beta 
\begin{pmatrix}
1  \\
z_1
\end{pmatrix}, 
\dfrac{1}{N} \begin{pmatrix}
1 & z_1^\top \\
z_1 & S_1 + z_1 z_1^\top
\end{pmatrix}\right)\\
z_1 | S_1 \sim \mathcal{N}\left(\mu_1, \dfrac{S_1}{\nu_0}\right)\\
S_1 \sim \mathcal{IW}(\nu_0 \bar S_1, \nu_0 + J + 1)\\
\beta \sim \mathcal{N}(0, \sigma_0^2)
\end{cases}
\end{align}

This model leads to
\[
\boxed{
p(x\mid\gamma=1)
\propto
|\bar R|^{-1/2}
\left(1+\frac{Nx_j^2}{\nu_0}\right)^{(\nu_0+3)/2}
\left(x^\top\bar R^{-1}x+\frac{\nu_0}{N}\right)^{-(\nu_0+J+2)/2}
\Normal\left(x_1\mid 0,\frac{1}{N} + \sigma_0^2\right)
}.
\]
Since $|\bar R|$ and $x^\top\bar R^{-1}x$ do not depend on which SNP is chosen,
the PIP weights simplify further. Dropping the common factors gives
\begin{equation}\label{eq:PIP-model-B}
\boxed{
p(\gamma=j \mid x)
\propto
\left(1+\frac{Nx_j^2}{\nu_0}\right)^{(\nu_0+3)/2}
\Normal\left(x_j \mid 0,\frac{1}{N} + \sigma_0^2\right)
}.
\end{equation}

\paragraph{Compare to SER-RSS's PIP.} Recall SER-RSS's PIP:
\begin{equation}\label{eq:PIP-susie-RSS}
\tilde{p}(\gamma=j\mid x)
\propto \exp\left(\dfrac{N\sigma_0^{2}}{N\sigma_0^{2} + 1} \dfrac{N x_j^2}{2} \right)
\end{equation}
We now show that~\eqref{eq:PIP-model-B} tends to~\eqref{eq:PIP-susie-RSS}
as $\nu_0\to\infty$. Since
\[
\left(1+\frac{Nx_j^2}{\nu_0}\right)^{(\nu_0+3)/2} = \left(1+ \dfrac{2}{\nu_0 + 3} \times \dfrac{\nu_0 + 3}{\nu_0} \frac{Nx_j^2}{2}\right)^{(\nu_0+3)/2} 
\asymp \exp\left\{\frac{\nu_0 + 3}{\nu_0} \dfrac{Nx_j^2}{2}\right\}
\]
Also,
\[
\Normal\left(x_j\mid 0,\frac{1}{N}+\sigma_0^2\right)
\propto
\exp\left\{
-\frac{x_j^2}{2(1/N+\sigma_0^2)}
\right\}
=
\exp\left\{
-\frac{1}{1+N\sigma_0^2} \frac{N x_j^2}{2}
\right\}.
\]
Therefore the limiting weight in~\eqref{eq:PIP-model-B} is asymptotically (as $\nu_0$ gets large)
\[
\exp\left\{\left(\frac{\nu_0 + 3}{\nu_0} - \dfrac{1}{N\sigma_0^2 + 1}\right) \dfrac{Nx_j^2}{2}\right\} = \exp\left\{\left(\dfrac{N\sigma_0^2}{N\sigma_0^2 + 1} - \dfrac{1}{\nu_0 + 3}\right) \dfrac{Nx_j^2}{2}\right\} \approx \exp\left\{\dfrac{N\sigma_0^2}{N\sigma_0^2 + 1} \dfrac{Nx_j^2}{2}\right\}
\]
which is exactly the SER-RSS PIP weight in~\eqref{eq:PIP-susie-RSS}. Hence, for finite $\nu_0$, PIP under SER-love model (B) is slightly more diffused than susie-rss. 

\paragraph{Inference under model B~\eqref{eq:SER-model-B}.}
For a general SNP $j$, write
\[
\mu_j=\bar R_{-j,j},
\qquad
\bar S_j=\bar R_{-j,-j}-\mu_j\mu_j^\top,
\qquad
\kappa_j=\nu_0+Nx_j^2.
\]
Under model B, conditional on $\gamma=j$, the likelihood factorizes as
\[
p(x\mid \beta,z_j,S_j,\gamma=j)
=
\Normal\left(x_j\mid \beta,\frac{1}{N}\right)
\Normal_{J-1}\left(x_{-j}\mid x_jz_j,\frac{S_j}{N}\right).
\]
Thus $\beta$ is conditionally independent of $(z_j,S_j)$ given $x$ and
$\gamma=j$, and
\[
p(\beta,z_j,S_j\mid x,\gamma=j)
=
p(\beta\mid x_j,\gamma=j)\,
p(z_j,S_j\mid x,\gamma=j).
\]

$\bullet$ The posterior for $\beta$ is the usual normal-normal update:
\[
\boxed{
\beta\mid x,\gamma=j
\sim
\Normal\left(
\frac{N}{N+\sigma_0^{-2}}x_j,\,
\frac{1}{N+\sigma_0^{-2}}
\right).
}
\]

$\bullet$ For the LD-column variables, first condition on $S_j$. Combining
\[
z_j\mid S_j\sim \Normal_{J-1}\left(\mu_j,\frac{S_j}{\nu_0}\right),
\qquad
x_{-j}\mid x_j,z_j,S_j,\gamma=j
\sim
\Normal_{J-1}\left(x_jz_j,\frac{S_j}{N}\right)
\]
gives
\[
\boxed{
z_j\mid S_j,x,\gamma=j
\sim
\Normal_{J-1}\left(
\frac{\nu_0\mu_j+Nx_jx_{-j}}{\nu_0+Nx_j^2},\,
\frac{S_j}{\nu_0+Nx_j^2}
\right).
}
\]

$\bullet$ Integrating out $z_j$ gives
\[
x_{-j}\mid x_j,S_j,\gamma=j
\sim
\Normal_{J-1}\left(
x_j\mu_j,\,
\left(\frac{1}{N}+\frac{x_j^2}{\nu_0}\right)S_j
\right).
\]
Therefore the posterior for $S_j$ is inverse-Wishart:
\[
\boxed{
S_j\mid x,\gamma=j
\sim
\mathcal{IW}\left(
\nu_0\bar S_j+
\frac{N\nu_0}{\nu_0+Nx_j^2}
(x_{-j}-x_j\mu_j)(x_{-j}-x_j\mu_j)^\top,\,
\nu_0+J+2
\right).
}
\]
Equivalently, the joint posterior of $(z_j,S_j)$ can be represented by first
drawing $S_j$ from this inverse-Wishart distribution and then drawing $z_j$
from the conditional normal distribution above. Finally, the full posterior
over $\gamma$ uses the PIP weights in~\eqref{eq:PIP-model-B}, normalized over
$j=1,\ldots,J$.
\end{document}
